\begin{proposition}[What the arity-one equations fix]
\label{prop:c11-arity1}
\begin{itemize}
\item (i) Along $\Delta$ the quadratic part of \eqref{eq:c9-arity1} fixes the translation-invariant part $\big[-\tfrac23\int(A,A^*)+\tfrac13\int(\lambda^*,\lambda)+\tfrac43\int(c^*,c)\big]$ units (Theorem \ref{thm:c9-quadratic} (iii)) and leaves $\theta$, $\theta_\lambda$, $\theta_V$, $\kappa_V$, $\kappa_\lambda$ free. Along $P_b$ the quadratic part of the arity-one obstruction vanishes (Proposition \ref{prop:c11-position} (i)), so the quadratic part of $\hbar F_1^{(1)}(P_b)$ is $Q$-closed and every coefficient of Definition \ref{def:c11-free} along $P_b$ is free in arity one, the arity-two equations relating them to the coefficients along $\Sigma$ (Theorem \ref{thm:c11-oddar2}); the translation-covariant choice, $\hbar F_1^{(1)}(P_b)$ equal to half the position part along $b$ of $\hbar F_1^{(1)}(\Delta)$, solves the equation along $P_b$ at every order in the fields.
\item (ii) Two solutions along $\Delta$ differ by an $s$-closed local functional of degree $-1$ and Weyl weight $4$. The common shift of $\theta$, $\theta_\lambda$, $\theta_V$ by $t$ adds $t$ units times $S^{(1)}(\Delta)$, which is $s$-closed (Lemma \ref{lem:c9-SDelta} (iii)); the directions $Q\int(\iota_xA|_V,c^*)$ and $Q\int(\rho(x)\lambda^*,\lambda^*)$ complete to $s\int(\iota_xA|_V,c^*)$ and $s\int(\rho(x)\lambda^*,\lambda^*)$, and an $s$-exact direction $sG$ contributes $s\{G,O\}$ to $D_1(O)$ on cocycles, zero on cohomology.
\item (iii) Since $\{I,S^{(1)}(\Delta)\}=0$ while the brackets of $I$ with the sector parts are nonzero, the cubic part of \eqref{eq:c9-arity1} depends on the sector coefficients through $(\theta-\theta_V)$ times the bracket of the cubic part of $I$ with the sector part of $A|_{W_{\mathbb R}^\perp}$ plus $(\theta_\lambda-\theta_V)$ times its bracket with the sector part of $\lambda$. It is solvable exactly when this sum is congruent to the cubic part of the composite part minus the bracket of the cubic part of $I$ with the translation-invariant part, modulo the images under $Q$ of the cubic local functionals of degree $-1$. If the classes of the two brackets modulo the images under $Q$ of the cubic local functionals of degree $-1$, Weyl weight $4$ and degree at most one in the position are linearly independent, the cubic equation fixes $\theta-\theta_V$ and $\theta_\lambda-\theta_V$.
\end{itemize}
\end{proposition}
